% GENERATED FILE. Edit canonical lessons, then run npm run generate:books.
% canonical-source-hash: fnv1a64:79d85893adf3c80a
% canonical-lessons: ES-C751-aficion, ES-C751-pasatiempo, ES-C751-coleccionar, ES-C751-coleccion, ES-C751-acuarela, ES-R751-repaso-744-751

\chapter{Mis aficiones}
\label{ch:es-mis-aficiones}
\hlchaptermodality{751}

\begin{chapteropening}
\textbf{By the end of this chapter:} \emph{I can talk about my hobbies, what I collect and what I like to paint.}

Complete the last lesson of chapter 751: I can talk about my hobbies, what I collect and what I like to paint.
\end{chapteropening}

\section[la afición]{la afición — the hobby; liking; fans}

\label{lesson:ES-C751-aficion}



\begin{quote}
Before the new one: say the Spanish for the colleague, then the Spanish for support.
\end{quote}

\subsection*{You'll want to know: la afición}
\textbf{la afición} — \textquotedblleft{}the hobby; liking; fans\textquotedblright{}.

\begin{quote}
Mi afición es la fotografía.
\end{quote}

\begin{quote}
\emph{My hobby is photography.}
\end{quote}

\emph{Tener afición a} — to be keen on. \emph{La afición} is also a team's fans.

\begin{cousinweb}
\emph{afición} comes from Latin \emph{affectiō} \textquotedblleft{}disposition, affection\textquotedblright{} (EtymDB).
\end{cousinweb}

\begin{grammarlens}[title={What you've built}]
One more word for telling, explaining and describing at B1.
\end{grammarlens}

\subsection*{Your turn}
\begin{itemize}
\raggedright
  \item \emph{Say it:} \emph{la afición}
  \item \emph{Say it:} \emph{la afición}, once more
  \item \emph{Say it:} a sentence of your own with \emph{la afición}
  \item \emph{Recall:} say the Spanish for the colleague, then the Spanish for support, then say \emph{la afición} again
\end{itemize}

\begin{tcolorbox}[breakable,colback=teal!4,colframe=teal!35!black,title={Before you move on}]
What does la afición mean? (\textquotedblleft{}The hobby; liking; fans\textquotedblright{}.) Say it once more.
\end{tcolorbox}
\section[el pasatiempo]{el pasatiempo — the pastime}

\label{lesson:ES-C751-pasatiempo}



\begin{quote}
Before the new one: say the Spanish for support, then the Spanish for the hobby.
\end{quote}

\subsection*{You'll want to know: el pasatiempo}
\textbf{el pasatiempo} — \textquotedblleft{}the pastime\textquotedblright{}.

\begin{quote}
Leer es mi pasatiempo favorito.
\end{quote}

\begin{quote}
\emph{Reading is my favourite pastime.}
\end{quote}

Plural \emph{los pasatiempos} — also puzzles in a newspaper.

\begin{cousinweb}
\emph{pasatiempo} is built on \emph{pasar} \textquotedblleft{}to pass\textquotedblright{} and \emph{tiempo} \textquotedblleft{}time\textquotedblright{}. \emph{Pasar} is usually traced to a spoken Latin verb built on \emph{passus} \textquotedblleft{}step\textquotedblright{}, and \emph{tiempo} comes from Latin \emph{tempus} \textquotedblleft{}time\textquotedblright{} (EtymDB).
\end{cousinweb}

\begin{grammarlens}[title={What you've built}]
One more word for telling, explaining and describing at B1.
\end{grammarlens}

\subsection*{Your turn}
\begin{itemize}
\raggedright
  \item \emph{Say it:} \emph{el pasatiempo}
  \item \emph{Say it:} \emph{el pasatiempo}, once more
  \item \emph{Say it:} a sentence of your own with \emph{el pasatiempo}
  \item \emph{Recall:} say the Spanish for support, then the Spanish for the hobby, then say \emph{el pasatiempo} again
\end{itemize}

\begin{tcolorbox}[breakable,colback=teal!4,colframe=teal!35!black,title={Before you move on}]
What does el pasatiempo mean? (\textquotedblleft{}The pastime\textquotedblright{}.) Say it once more.
\end{tcolorbox}
\section[coleccionar]{coleccionar — to collect}

\label{lesson:ES-C751-coleccionar}



\begin{quote}
Before the new one: say the Spanish for the hobby, then the Spanish for the pastime.
\end{quote}

\subsection*{You'll want to know: coleccionar}
\textbf{coleccionar} — \textquotedblleft{}to collect\textquotedblright{}.

\begin{quote}
Colecciono sellos desde niño.
\end{quote}

\begin{quote}
\emph{I have collected stamps since I was a boy.}
\end{quote}

For hobbies; to pick someone up is \emph{recoger}.

\begin{cousinweb}
\emph{coleccionar} is built on \emph{colección}, a learned borrowing from Latin \emph{collēctiō} \textquotedblleft{}a gathering together\textquotedblright{} (EtymDB).
\end{cousinweb}

\begin{grammarlens}[title={What you've built}]
One more word for telling, explaining and describing at B1.
\end{grammarlens}

\subsection*{Your turn}
\begin{itemize}
\raggedright
  \item \emph{Say it:} \emph{coleccionar}
  \item \emph{Say it:} \emph{coleccionar}, once more
  \item \emph{Say it:} a sentence of your own with \emph{coleccionar}
  \item \emph{Recall:} say the Spanish for the hobby, then the Spanish for the pastime, then say \emph{coleccionar} again
\end{itemize}

\begin{tcolorbox}[breakable,colback=teal!4,colframe=teal!35!black,title={Before you move on}]
What does coleccionar mean? (\textquotedblleft{}To collect\textquotedblright{}.) Say it once more.
\end{tcolorbox}
\section[la colección]{la colección — the collection}

\label{lesson:ES-C751-coleccion}



\begin{quote}
Before the new one: say the Spanish for the pastime, then the Spanish for to collect.
\end{quote}

\subsection*{You'll want to know: la colección}
\textbf{la colección} — \textquotedblleft{}the collection\textquotedblright{}.

\begin{quote}
Tiene una colección de monedas antiguas.
\end{quote}

\begin{quote}
\emph{He has a collection of old coins.}
\end{quote}

\begin{cousinweb}
\emph{colección} was borrowed from Latin \emph{collēctiō} \textquotedblleft{}a gathering together\textquotedblright{} (EtymDB), from \emph{colligere}, built on \emph{legere} \textquotedblleft{}to gather\textquotedblright{}.
\end{cousinweb}

\begin{grammarlens}[title={What you've built}]
One more word for telling, explaining and describing at B1.
\end{grammarlens}

\subsection*{Your turn}
\begin{itemize}
\raggedright
  \item \emph{Say it:} \emph{la colección}
  \item \emph{Say it:} \emph{la colección}, once more
  \item \emph{Say it:} a sentence of your own with \emph{la colección}
  \item \emph{Recall:} say the Spanish for the pastime, then the Spanish for to collect, then say \emph{la colección} again
\end{itemize}

\begin{tcolorbox}[breakable,colback=teal!4,colframe=teal!35!black,title={Before you move on}]
What does la colección mean? (\textquotedblleft{}The collection\textquotedblright{}.) Say it once more.
\end{tcolorbox}
\section[la acuarela]{la acuarela — watercolour}

\label{lesson:ES-C751-acuarela}



\begin{quote}
Before the new one: say the Spanish for to collect, then the Spanish for the collection.
\end{quote}

\subsection*{You'll want to know: la acuarela}
\textbf{la acuarela} — \textquotedblleft{}watercolour\textquotedblright{}.

\begin{quote}
Pinto acuarelas los fines de semana.
\end{quote}

\begin{quote}
\emph{I paint watercolours at weekends.}
\end{quote}

Both the paint and the picture.

\begin{cousinweb}
\emph{acuarela} was borrowed from Italian \emph{acquerello} (EtymDB), built on \emph{acqua} \textquotedblleft{}water\textquotedblright{}, from Latin \emph{aqua}.
\end{cousinweb}

\begin{grammarlens}[title={What you've built}]
One more word for telling, explaining and describing at B1.
\end{grammarlens}

\subsection*{Your turn}
\begin{itemize}
\raggedright
  \item \emph{Say it:} \emph{la acuarela}
  \item \emph{Say it:} \emph{la acuarela}, once more
  \item \emph{Say it:} a sentence of your own with \emph{la acuarela}
  \item \emph{Recall:} say the Spanish for to collect, then the Spanish for the collection, then say \emph{la acuarela} again
\end{itemize}

\begin{tcolorbox}[breakable,colback=teal!4,colframe=teal!35!black,title={Before you move on}]
What does la acuarela mean? (\textquotedblleft{}Watercolour\textquotedblright{}.) Say it once more.
\end{tcolorbox}
\section[(dialogue)]{Repaso de los capítulos 744–751}

\label{lesson:ES-R751-repaso-744-751}



\begin{quote}
Say the Spanish for love at first sight and the compliment.
\end{quote}

\subsection*{Your turn}
Say these aloud:

\begin{itemize}
\raggedright
  \item love at first sight, the compliment, the courtship, the suitor, closeness
  \item the relationship, the break-up, unfaithful, to betray, betrayal
  \item to stroke, the deception, to chat up, to flirt, the lover
  \item the brother-in-law, the father-in-law, the mother-in-law, the son-in-law, the daughter-in-law
  \item to adopt, to take in, adoption, the stepmother, the stepfather
  \item the twin, the stepbrother, the teenager, the great-grandfather, the orphan
  \item the bond, to spoil, to live together, the colleague, support
  \item the hobby, the pastime, to collect, the collection, watercolour
\end{itemize}

\begin{tcolorbox}[breakable,colback=teal!4,colframe=teal!35!black,title={Before you move on}]
Love at first sight? (\textbf{el flechazo}.) The compliment? (\textbf{el piropo}.) The courtship? (\textbf{el noviazgo}.) The suitor? (\textbf{el pretendiente}.) Closeness? (\textbf{la complicidad}.) The relationship? (\textbf{la relación}.) The break-up? (\textbf{la ruptura}.) Unfaithful? (\textbf{infiel}.) To betray? (\textbf{traicionar}.) Betrayal? (\textbf{la traición}.) To stroke? (\textbf{acariciar}.) The deception? (\textbf{el engaño}.) To chat up? (\textbf{ligar}.) To flirt? (\textbf{coquetear}.) The lover? (\textbf{el amante}.) The brother-in-law? (\textbf{el cuñado}.) The father-in-law? (\textbf{el suegro}.) The mother-in-law? (\textbf{la suegra}.) The son-in-law? (\textbf{el yerno}.) The daughter-in-law? (\textbf{la nuera}.) To adopt? (\textbf{adoptar}.) To take in? (\textbf{acoger}.) Adoption? (\textbf{la adopción}.) The stepmother? (\textbf{la madrastra}.) The stepfather? (\textbf{el padrastro}.) The twin? (\textbf{el gemelo}.) The stepbrother? (\textbf{el hermanastro}.) The teenager? (\textbf{el adolescente}.) The great-grandfather? (\textbf{el bisabuelo}.) The orphan? (\textbf{el huérfano}.) The bond? (\textbf{el vínculo}.) To spoil? (\textbf{mimar}.) To live together? (\textbf{convivir}.) The colleague? (\textbf{el colega}.) Support? (\textbf{el apoyo}.) The hobby? (\textbf{la afición}.) The pastime? (\textbf{el pasatiempo}.) To collect? (\textbf{coleccionar}.) The collection? (\textbf{la colección}.) Watercolour? (\textbf{la acuarela}.)
\end{tcolorbox}
